\documentclass[11pt]{article}

\usepackage[utf8]{inputenc}
\usepackage[T1]{fontenc}
\usepackage[spanish]{babel}
\usepackage[a4paper,landscape,left=0cm,right=0cm,top=0.2cm,bottom=0cm]{geometry}
\usepackage{tikz}
\usepackage{xcolor}
\usepackage{graphicx}
\usepackage{lmodern}
\pagestyle{empty}
\setlength{\parindent}{0pt}

\definecolor{azulUVE}{RGB}{0,120,180}
\definecolor{naranjaUVE}{RGB}{220,90,20}

\begin{document}
\thispagestyle{empty}
\centering

\resizebox{\paperwidth}{!}{%
\begin{tikzpicture}[x=1cm,y=1cm]

% Caja total del dibujo
\path[use as bounding box] (0,0) rectangle (27,18);

%------------------ TITULO ------------------
\node[align=center] at (13.5,17.1) {
    {\LARGE\bfseries\textcolor{azulUVE}{V de Gowin para estadística}}\\[0.18cm]
    {\normalsize\bfseries Tema: Análisis de la relación entre el uso de las redes sociales y el rendimiento académico en estudiantes universitarios}
};

%------------------ ENCABEZADOS ------------------
\node[anchor=west, align=left] at (0.7,15.6) {
    {\large\bfseries\textcolor{azulUVE}{PARTE CONCEPTUAL}}\\[0.05cm]
    \textbf{Pienso sobre el tema}
};

\node[anchor=west, align=left] at (19.6,15.6) {
    {\large\bfseries\textcolor{azulUVE}{PARTE METODOLÓGICA}}\\[0.05cm]
    \textbf{Actúo sobre el tema}
};

%------------------ LINEAS DE LA V ------------------
\draw[line width=1pt] (7.6,14.6) -- (13.5,3.0);
\draw[line width=1pt] (19.4,14.6) -- (13.5,3.0);

%------------------ 1 PREGUNTA ------------------
\node[
    text width=6.6cm,
    align=center,
    font=\small
] at (13.5,13.8) {
    {\fontsize{22}{22}\selectfont\bfseries\textcolor{azulUVE}{1}}\\
    {\large\bfseries Pregunta clave}\\[0.04cm]
    {\color{naranjaUVE}\bfseries ¿Qué queremos conocer?}\\[0.16cm]
    ¿Existe una relación entre el tiempo dedicado a las redes sociales y el rendimiento académico en estudiantes universitarios?
};

%------------------ 2 OBJETO ------------------
\node[
    text width=8.7cm,
    align=center,
    font=\small
] at (13.5,1.95) {
    {\fontsize{22}{22}\selectfont\bfseries\textcolor{azulUVE}{2}}\\
    {\large\bfseries Objeto de estudio}\\[0.04cm]
    {\color{naranjaUVE}\bfseries ¿Qué fenómeno estudiamos?}\\[0.16cm]
    La relación entre el uso de redes sociales, especialmente el tiempo de uso, y el rendimiento académico en estudiantes universitarios, considerando su posible influencia sobre la concentración.
};

%------------------ LADO IZQUIERDO ------------------

% 3 Conceptos
\node[
    anchor=north west,
    text width=7.15cm,
    align=left,
    font=\fontsize{8.7}{10.4}\selectfont
] at (0.7,15) {
    {\fontsize{22}{22}\selectfont\bfseries\textcolor{azulUVE}{3}} {\large\bfseries Conceptos relacionados}\\
    {\color{naranjaUVE}\bfseries ¿Cuáles son los conceptos clave?}\\[0.16cm]
    \textbf{Redes sociales:} plataformas digitales que permiten interacción, creación y difusión de contenido entre usuarios (Kaplan y Haenlein, 2010).\\[0.14cm]
    \textbf{Rendimiento académico:} nivel de logro alcanzado por un estudiante, medido mediante notas, promedio u otros indicadores educativos (Chadwick, 1979).\\[0.14cm]
    \textbf{Concentración:} capacidad de mantener la atención sostenida y reducir distracciones durante el aprendizaje (APA, s.f.).
};

% 5 Principios
\node[
    anchor=north west,
    text width=8.15cm,
    align=left,
    font=\fontsize{8.7}{10.4}\selectfont
] at (0.7,9.6) {
    {\fontsize{22}{22}\selectfont\bfseries\textcolor{azulUVE}{5}} {\large\bfseries Principios}\\
    {\color{naranjaUVE}\bfseries ¿Cómo sucede el fenómeno?}\\[0.16cm]
    \textbf{Relación correlacional entre variables:} dos variables pueden variar de manera asociada sin que esto implique causalidad directa.\\[0.14cm]
    \textbf{Multifactorialidad del rendimiento académico:} el desempeño estudiantil depende de factores como hábitos de estudio, estrés, motivación, contexto y uso de tecnología.
};

% 7 Teorias
\node[
    anchor=north west,
    text width=6.15cm,
    align=left,
    font=\fontsize{8.7}{10.4}\selectfont
] at (0.7,5.7) {
    {\fontsize{22}{22}\selectfont\bfseries\textcolor{azulUVE}{7}} {\large\bfseries Teorías}\\
    {\color{naranjaUVE}\bfseries ¿Por qué sucede?}\\[0.16cm]
    \textbf{Teoría del desplazamiento del tiempo:} el tiempo dedicado a redes sociales puede restar horas al estudio, descanso o actividades académicas.\\[0.14cm]
    \textbf{Teoría de la distracción cognitiva / atención limitada:} las notificaciones, interrupciones y estímulos digitales pueden reducir la concentración y afectar el aprendizaje.
};

%------------------ LADO DERECHO ------------------

% 4 Procedimiento
\node[
    anchor=north west,
    text width=7.15cm,
    align=left,
    font=\fontsize{8.7}{10.4}\selectfont
] at (19.6,15) {
    {\fontsize{22}{22}\selectfont\bfseries\textcolor{azulUVE}{4}} {\large\bfseries Procedimiento}\\
    {\color{naranjaUVE}\bfseries ¿Qué haremos para observar el fenómeno?}\\[0.16cm]
    1. Identificar las variables relevantes del dataset.\\[0.06cm]
    2. Limpiar y organizar los datos.\\[0.06cm]
    3. Realizar análisis descriptivo.\\[0.06cm]
    4. Aplicar análisis correlacional entre uso de redes y rendimiento.\\[0.06cm]
    5. Explorar el papel de la concentración u otras variables asociadas.\\[0.06cm]
    6. Interpretar los resultados con apoyo bibliográfico.
};

% 6 Registro
\node[
    anchor=north west,
    text width=7.15cm,
    align=left,
    font=\fontsize{8.7}{10.4}\selectfont
] at (19.6,9.7) {
    {\fontsize{22}{22}\selectfont\bfseries\textcolor{azulUVE}{6}} {\large\bfseries Registro y transformación de datos}\\
    {\color{naranjaUVE}\bfseries ¿Qué medimos directamente?}\\[0.16cm]
    \textbf{Fuente:} Student Social Media, Academic Performance Dataset (Kaggle).\\[0.08cm]
    \textbf{Unidad de análisis:} estudiantes universitarios.\\[0.08cm]
    \textbf{Variable independiente:} tiempo dedicado a redes sociales.\\[0.08cm]
    \textbf{Variable dependiente:} rendimiento académico.\\[0.08cm]
    \textbf{Otras variables:} concentración, hábitos de estudio, estrés u otras disponibles.\\[0.08cm]
    \textbf{Transformaciones:} depuración, recodificación, tablas, gráficos y medidas de asociación.
};

% 8 Conclusiones
\node[
    anchor=north west,
    text width=7.15cm,
    align=left,
    font=\fontsize{8.7}{10.4}\selectfont
] at (19.6,3.9) {
    {\fontsize{22}{22}\selectfont\bfseries\textcolor{azulUVE}{8}} {\large\bfseries Conclusiones preliminares}\\
    {\color{naranjaUVE}\bfseries ¿Qué podemos afirmar?}\\[0.16cm]
    1. Puede existir una relación entre uso de redes sociales y rendimiento académico.\\[0.06cm]
    2. Esa relación podría no ser fuerte ni lineal.\\[0.06cm]
    3. La concentración puede funcionar como variable mediadora.\\[0.06cm]
    4. La literatura reporta resultados mixtos, por lo que el análisis empírico será decisivo.
};

%------------------ REFERENCIAS ------------------
\node[
    anchor=west,
    text width=17 cm,
    align=left,
    font=\footnotesize
] at (0.7,0.00001) {
\textbf{Referencias:}
Kaplan y Haenlein (2010); Chadwick (1979); APA (s.f.); León Luyo et al. (2023); Santos (2010); Aljemely (2020); Ramírez Jaime et al. (2023); Rodelo Moreno y Lizárraga Reyes (2018); Miranda-Rosero et al. (2023).
};

\end{tikzpicture}%
}

\end{document}